\documentclass[11pt,a4paper]{article}
\usepackage{fontspec}
\usepackage{polyglossia}
\setmainlanguage{russian}
\setotherlanguage{english}
\setmainfont{Liberation Serif}
\usepackage[left=1.0cm,right=1.0cm,top=0.8cm,bottom=0.7cm]{geometry}
\usepackage{amsmath,amssymb}
\usepackage{array}
\usepackage{multirow}
\setlength{\parindent}{0pt}
\pagestyle{empty}

\newcounter{zad}


\newlength{\otvsh}\setlength{\otvsh}{3.5cm}
\newlength{\plsh}\setlength{\plsh}{0.7cm}
\newlength{\polesh}\setlength{\polesh}{4.6cm}
\newlength{\uslsh}\setlength{\uslsh}{13.6cm}
\newlength{\podnyat}\setlength{\podnyat}{1.1ex}
\newlength{\minpole}\setlength{\minpole}{1.3cm}

\newsavebox{\uslbox}
\newlength{\uslvys}
\newlength{\blh}

\newcommand{\calcblh}[1]{%
  \setlength{\blh}{\dimexpr(\uslvys-\arrayrulewidth)/#1-\arrayrulewidth\relax}}
\newcommand{\mt}[1]{\parbox[t][\blh][t]{\otvsh}{\small\bfseries #1\strut}}

\newcommand{\poleA}{\calcblh{1}\setlength{\tabcolsep}{3pt}%
  \begin{tabular}{|p{\otvsh}|p{\plsh}|}\hline \mt{} & \\ \hline\end{tabular}}
\newcommand{\poleB}[2]{\calcblh{2}\setlength{\tabcolsep}{3pt}%
  \begin{tabular}{|p{\otvsh}|p{\plsh}|}\hline
  \mt{#1} & \multirow{2}{*}{} \\ \cline{1-1} \mt{#2} & \\ \hline\end{tabular}}
\newcommand{\poleC}[3]{\calcblh{3}\setlength{\tabcolsep}{3pt}%
  \begin{tabular}{|p{\otvsh}|p{\plsh}|}\hline
  \mt{#1} & \multirow{3}{*}{} \\ \cline{1-1} \mt{#2} & \\ \cline{1-1}
  \mt{#3} & \\ \hline\end{tabular}}
\newcommand{\poleD}[4]{\calcblh{4}\setlength{\tabcolsep}{3pt}%
  \begin{tabular}{|p{\otvsh}|p{\plsh}|}\hline
  \mt{#1} & \multirow{4}{*}{} \\ \cline{1-1} \mt{#2} & \\ \cline{1-1}
  \mt{#3} & \\ \cline{1-1} \mt{#4} & \\ \hline\end{tabular}}

\newlength{\mezhzad}\setlength{\mezhzad}{14pt}
\newcommand{\znak}[1]{\ifx&#1&\else$#1$\ \fi}

\newcommand{\zad}[3]{%
  \stepcounter{zad}%
  \sbox{\uslbox}{\parbox[t]{\uslsh}{\strut\textbf{\thezad\ifx&#1&\else$^{*}$\fi.}\ #3\strut}}%
  \setlength{\uslvys}{\dimexpr\ht\uslbox+\dp\uslbox\relax}%
  \ifdim\uslvys<\minpole\setlength{\uslvys}{\minpole}\fi
  \par\vspace{\mezhzad}\noindent
  \begin{minipage}[t]{\polesh}\vspace{0pt}\vspace*{-\podnyat}#2\end{minipage}%
  \hfill
  \begin{minipage}[t]{\uslsh}\vspace{0pt}\usebox{\uslbox}\end{minipage}\par}


% ЛИСТОК КРУЖКА 7И, 29.09 — «Построй или докажи, что нельзя» (конструктивы). v3.
% Макросы — из listok_23-09_logika.tex; звёздочка — верхним индексом у номера (как в listok_15-09_logika.tex, \zadz).
% Скелет по решениям владельца 29.09 (00:40, 02:04, 02:36): три мягких входа (1–3), четыре главных (4–7),
% две звёздочки (8, 9), десятая — для сильного школьника, одна звёздочка, пункты-ступени.
% v2: 6 → правильный треугольник; 7 → «все наборы различных»; в 2 «хотя бы двух» (KRITIKA-29-09.md).
% v3 (POISK-cpm-mmm.md, POISK-vershina.md): 3 → ряд из 10/9 чисел (Малый мехмат 2024);
%     8* → «девять задач, по шесть отметок» (ЦПМ, Хамовники, «Мнимые противоречия», 2024) + пункт «по 7»;
%     10* → «запутанный турнир» (проблема Шютте, k=2): вершина, в которой сходятся все линии.
% ЛИНИИ:
%   строим:                       1 → 5 (шаг +3) → 9* → 10(в,г) (шаг +1)
%   запрет одной величиной:       3(а) пары соседей → 4(а) концы проводов → 8*(б) отметки за задачу → 10(б) поражения
%   мнимое противоречие:          2 → 3(б) → 8*(а)
%   при каких n / найдите все:    4(б), 5, 6, 7 → 10(г)
% Проверка перебором: proverka_29-09_konstruktivy.py; верификатор: VERIFIKACIYA-29-09.md
\newcommand{\shapka}{%
  \noindent{\LARGE\bfseries Построй или докажи, что нельзя}%
  \hfill{\small Кружок, 7\,И\quad·\quad29 сентября}%
  \vspace{4pt}\hrule\vspace{8pt}
  \noindent Фамилия, имя:\ \hrulefill\rule[-3mm]{0pt}{8mm}\par
  \vspace{2pt}
  \setcounter{zad}{0}}

\setlength{\mezhzad}{10pt}

\begin{document}
\shapka

\zad{}{\poleC{а}{б}{в}}{Разрежьте квадрат на \textbf{(а)}~7 квадратов; \textbf{(б)}~6 квадратов;
\textbf{(в)}~8 квадратов. Квадраты не обязаны быть одинаковыми.}

\zad{}{\poleA}{Могут ли и сумма, и произведение нескольких (хотя бы двух) натуральных чисел
равняться 99?}

\zad{}{\poleB{а}{б}}{В ряд записаны \textbf{(а)}~10 чисел; \textbf{(б)}~9 чисел. Сумма любых двух соседних
чисел положительна. Может ли сумма всех чисел быть отрицательной?}

\zad{}{\poleB{а}{б}}{Телефоны соединяют проводами, каждый провод соединяет два телефона,
и никакие два телефона не соединены дважды.\newline
\textbf{(а)}~Можно ли соединить 15 телефонов так, чтобы каждый был соединён ровно с тремя другими?\newline
\textbf{(б)}~При каких $n$ так можно соединить $n$ телефонов?}

\zad{}{\poleC{а}{б}{в}}{\textbf{(а)}~Разрежьте правильный треугольник на 6, на 7 и на 8 правильных треугольников
(не обязательно одинаковых).\newline
\textbf{(б)}~Докажите, что его можно разрезать на $n$ правильных треугольников при любом $n\ge6$.\newline
\textbf{(в)}~Докажите, что на 2 и на 3 правильных треугольника его разрезать нельзя.}

\zad{}{\poleB{а}{б}}{На острове живут рыцари, которые всегда говорят правду, и лжецы,
которые всегда лгут. За круглым столом сидят островитяне, и каждый сказал:
\textit{«оба моих соседа~--- лжецы»}. Сколько рыцарей может сидеть за столом,
если всего островитян \textbf{(а)}~10; \textbf{(б)}~7? Найдите все варианты.}

\zad{}{\poleB{а}{б}}{\textbf{(а)}~При каких $N$ и сумма, и произведение нескольких (хотя бы двух)
натуральных чисел могут равняться $N$?\newline
\textbf{(б)}~Найдите все наборы из нескольких (хотя бы двух) \textit{различных} натуральных чисел,
у которых сумма равна произведению.}

\zad{\star}{\poleB{а}{б}}{На кружке было 9 задач, и каждую из них рассказали преподавателю Аня, Боря и Вера,
каждый в свой момент. Аня отмечает у себя задачу, если рассказала её раньше Бори; Боря~--- если рассказал раньше Веры;
Вера~--- если рассказала раньше Ани. Может ли у каждого из троих оказаться отмечено \textbf{(а)}~по 6 задач; \textbf{(б)}~по 7 задач?}

\zad{\star}{\poleA}{Докажите, что квадрат нельзя разрезать на 5 квадратов.}

\zad{\star}{\poleD{а}{б}{в}{г}}{В турнире не меньше двух участников, каждый сыграл с каждым по одному разу, ничьих не было. Назовём турнир
\textit{запутанным}, если для любых двух участников найдётся третий, который выиграл у обоих.\newline
\textbf{(а)}~Докажите, что в запутанном турнире каждый участник проиграл хотя бы трём.\newline
\textbf{(б)}~Докажите, что в запутанном турнире не меньше 7 участников.\newline
\textbf{(в)}~Придумайте запутанный турнир с 7 участниками.\newline
\textbf{(г)}~При каких $n$ существует запутанный турнир с $n$ участниками?}

\end{document}
